% GENERATED FILE. Edit canonical lessons, then run npm run generate:books.
% canonical-source-hash: fnv1a64:19cb7e631fc99533
% canonical-lessons: ZH-C301-tao, ZH-C301-huang, ZH-C301-fu-help, ZH-C301-guzhe, ZH-C301-guoran, ZH-R301-review-296-301

\chapter{Getting Out}
\label{ch:zh-getting-out}
\hlchaptermodality{301}

\begin{chapteropening}
\textbf{By the end of this chapter:} \emph{I can say people ran from danger, someone helped a person up, and it all went as expected.}

Complete the last lesson of chapter 301: I can say people ran from danger, someone helped a person up, and it all went as expected.
\end{chapteropening}

\section[táo]{táo — to run away; to flee}

\label{lesson:ZH-C301-tao}



\begin{quote}
Before the new one: say the Mandarin for the cause of death, then the Mandarin for to blow up.
\end{quote}

\subsection*{You'll want to know: táo}
\emph{táo} — \textquotedblleft{}to run away; to flee\textquotedblright{}.

\begin{quote}
\emph{fángzi zháohuǒ le, dàjiā kuài táo} — The house is on fire, everyone run!
\end{quote}

Its characters come later; for now it is a word for the ear.

\begin{cousinweb}
\emph{táo} is inherited from Middle Chinese \emph{daw} (Guangyun, in Baxter's transcription), which glosses it 'to go away', 'to avoid' and 'to flee'.
\end{cousinweb}

\begin{grammarlens}[title={What you've built}]
One more word for telling what happened.
\end{grammarlens}

\subsection*{Your turn}
\begin{itemize}
\raggedright
  \item \emph{Say it:} \emph{táo}
  \item \emph{Say it:} \emph{táo}, once more
  \item \emph{Say it:} \emph{fángzi zháohuǒ le, dàjiā kuài táo} — The house is on fire, everyone run!
  \item \emph{Recall:} say the Mandarin for the cause of death, then the Mandarin for to blow up, then say \emph{táo} again
\end{itemize}

\begin{tcolorbox}[breakable,colback=teal!4,colframe=teal!35!black,title={Before you move on}]
What does táo mean? (\textquotedblleft{}To run away; to flee\textquotedblright{}.) Say it once more.
\end{tcolorbox}
\section[huāng]{huāng — flustered; in a panic}

\label{lesson:ZH-C301-huang}



\begin{quote}
Before the new one: say the Mandarin for to blow up, then the Mandarin for to flee from danger.
\end{quote}

\subsection*{You'll want to know: huāng}
\emph{huāng} — \textquotedblleft{}flustered; in a panic\textquotedblright{}.

\begin{quote}
\emph{bié huāng, wǒmen xiān bàojǐng} — Don't panic, let's call the police first.
\end{quote}

Its characters come later; for now it is a word for the ear.

\begin{cousinweb}
\emph{huāng} is inherited from Middle Chinese \emph{xwangX} (Guangyun, in Baxter's transcription), which lists it in a two-syllable word meaning 'dazed'. The level tone of huāng is irregular for \emph{xwangX}.
\end{cousinweb}

\begin{grammarlens}[title={What you've built}]
One more word for telling what happened.
\end{grammarlens}

\subsection*{Your turn}
\begin{itemize}
\raggedright
  \item \emph{Say it:} \emph{huāng}
  \item \emph{Say it:} \emph{huāng}, once more
  \item \emph{Say it:} \emph{bié huāng, wǒmen xiān bàojǐng} — Don't panic, let's call the police first.
  \item \emph{Recall:} say the Mandarin for to blow up, then the Mandarin for to flee from danger, then say \emph{huāng} again
\end{itemize}

\begin{tcolorbox}[breakable,colback=teal!4,colframe=teal!35!black,title={Before you move on}]
What does huāng mean? (\textquotedblleft{}Flustered; in a panic\textquotedblright{}.) Say it once more.
\end{tcolorbox}
\section[fú]{fú — to help up; to hold up with a hand}

\label{lesson:ZH-C301-fu-help}



\begin{quote}
Before the new one: say the Mandarin for to flee from danger, then the Mandarin for in a panic.
\end{quote}

\subsection*{You'll want to know: fú}
\emph{fú} — \textquotedblleft{}to help up; to hold up with a hand\textquotedblright{}.

\begin{quote}
\emph{wǒ fú tā qǐlai} — I helped him up.
\end{quote}

Its characters come later; for now it is a word for the ear.

\begin{cousinweb}
\emph{fú} is inherited from Middle Chinese \emph{bju} (Guangyun, in Baxter's transcription), which glosses it 'to hold up' and 'to help'.
\end{cousinweb}

\begin{grammarlens}[title={What you've built}]
One more word for telling what happened.
\end{grammarlens}

\subsection*{Your turn}
\begin{itemize}
\raggedright
  \item \emph{Say it:} \emph{fú}
  \item \emph{Say it:} \emph{fú}, once more
  \item \emph{Say it:} \emph{wǒ fú tā qǐlai} — I helped him up.
  \item \emph{Recall:} say the Mandarin for to flee from danger, then the Mandarin for in a panic, then say \emph{fú} again
\end{itemize}

\begin{tcolorbox}[breakable,colback=teal!4,colframe=teal!35!black,title={Before you move on}]
What does fú mean? (\textquotedblleft{}To help up; to hold up with a hand\textquotedblright{}.) Say it once more.
\end{tcolorbox}
\section[gǔzhé]{gǔzhé — to break a bone}

\label{lesson:ZH-C301-guzhe}



\begin{quote}
Before the new one: say the Mandarin for in a panic, then the Mandarin for to help up someone who fell.
\end{quote}

\subsection*{You'll want to know: gǔzhé}
\emph{gǔzhé} — \textquotedblleft{}to break a bone\textquotedblright{}.

\begin{quote}
\emph{tā de tuǐ gǔzhé le} — He broke his leg.
\end{quote}

Its characters come later; for now it is a word for the ear.

\begin{cousinweb}
\emph{gǔzhé} is built from \emph{gǔ} 'a bone' and \emph{zhé} 'to snap', both inherited from Middle Chinese, \emph{kwot} and \emph{tsyet} (Guangyun, in Baxter's transcription), which glosses the first, quoting an older dictionary, 'the core of the flesh' and the second 'to bend and break'.
\end{cousinweb}

\begin{grammarlens}[title={What you've built}]
One more word for telling what happened.
\end{grammarlens}

\subsection*{Your turn}
\begin{itemize}
\raggedright
  \item \emph{Say it:} \emph{gǔzhé}
  \item \emph{Say it:} \emph{gǔzhé}, once more
  \item \emph{Say it:} \emph{tā de tuǐ gǔzhé le} — He broke his leg.
  \item \emph{Recall:} say the Mandarin for in a panic, then the Mandarin for to help up someone who fell, then say \emph{gǔzhé} again
\end{itemize}

\begin{tcolorbox}[breakable,colback=teal!4,colframe=teal!35!black,title={Before you move on}]
What does gǔzhé mean? (\textquotedblleft{}To break a bone\textquotedblright{}.) Say it once more.
\end{tcolorbox}
\section[\texorpdfstring{\zh{果然} (guǒrán)}{果然 (guǒrán)}]{\zh{果然} (guǒrán) — sure enough; as expected}

\label{lesson:ZH-C301-guoran}



\begin{quote}
Before the new one: say the Mandarin for to help up someone who fell, then the Mandarin for to break a bone.
\end{quote}

\subsection*{You'll want to know: \zh{果然}}
\textbf{\zh{果然}} — \emph{guǒrán} — \textquotedblleft{}sure enough; as expected\textquotedblright{}.

\begin{quote}
\textbf{\zh{他说会下雨果然下雨了}} \emph{tā shuō huì xiàyǔ, guǒrán xiàyǔ le} — He said it would rain, and sure enough it did.
\end{quote}

\begin{cousinweb}
\textbf{\zh{果然}} (\emph{guǒrán}) is built from \zh{果} 'resolute' and \zh{然} 'so', both inherited from Middle Chinese, \emph{kwaX} and \emph{nyen} (Guangyun, in Baxter's transcription), which glosses the first 'resolute' and 'certain' and the second 'so, like this': so, for certain.
\end{cousinweb}

\begin{grammarlens}[title={What you've built}]
One more word for telling what happened.
\end{grammarlens}

\subsection*{Your turn}
\begin{itemize}
\raggedright
  \item \emph{Say it:} \emph{guǒrán}
  \item \emph{Say it:} \emph{guǒrán}, once more
  \item \emph{Say it:} \emph{tā shuō huì xiàyǔ, guǒrán xiàyǔ le} — He said it would rain, and sure enough it did.
  \item \emph{Recall:} say the Mandarin for to help up someone who fell, then the Mandarin for to break a bone, then say \emph{guǒrán} again
\end{itemize}

\begin{tcolorbox}[breakable,colback=teal!4,colframe=teal!35!black,title={Before you move on}]
What does \zh{果然} mean? (\textquotedblleft{}Sure enough; as expected\textquotedblright{}.) Say it once more.
\end{tcolorbox}
\section[(dialogue)]{Looking back: from A Crash to Getting Out}

\label{lesson:ZH-R301-review-296-301}



\begin{quote}
Say the Mandarin for something unexpected and a car crash.
\end{quote}

\subsection*{Your turn}
\begin{itemize}
\raggedright
  \item \emph{Say it:} something unexpected, a car crash, to put on the brakes, a beginner at the wheel, to pay for damage
  \item \emph{Say it:} insurance for the car, a serious accident, who is to blame, what comes of it, to deal with a problem
  \item \emph{Say it:} a fire that destroys buildings, to put out a fire, a firefighter, to shout loudly, the cry for help
  \item \emph{Say it:} to be saved from danger, before it was too late, first aid, to pass out, to fight to save a patient
  \item \emph{Say it:} an earthquake, to die, the life of a living thing, the cause of death, to blow up
  \item \emph{Say it:} to flee from danger, in a panic, to help up someone who fell, to break a bone, sure enough
  \item \emph{Read it:} \zh{故} \zh{救} \zh{死} — each one aloud, then a word that holds it
\end{itemize}

\begin{tcolorbox}[breakable,colback=teal!4,colframe=teal!35!black,title={Before you move on}]
Something unexpected? (\textbf{\zh{意外}}, \emph{yìwài}.) A car crash? (\emph{chēhuò}.) To put on the brakes? (\emph{shāchē}.) A beginner at the wheel? (\textbf{\zh{新手}}, \emph{xīnshǒu}.) To pay for damage? (\emph{péiqián}.) Insurance for the car? (\emph{bǎoxiǎn}.) A serious accident? (\textbf{\zh{事故}}, \emph{shìgù}.) Who is to blame? (\emph{zérèn}.) What comes of it? (\textbf{\zh{后果}}, \emph{hòuguǒ}.) To deal with a problem? (\emph{chǔlǐ}.) A fire that destroys buildings? (\emph{huǒzāi}.) To put out a fire? (\textbf{\zh{救火}}, \emph{jiùhuǒ}.) A firefighter? (\emph{xiāofángyuán}.) To shout loudly? (\emph{hǎn}.) The cry for help? (\emph{jiùmìng}.) To be saved from danger? (\textbf{\zh{得救}}, \emph{déjiù}.) Before it was too late? (\emph{jíshí}.) First aid? (\emph{jíjiù}.) To pass out? (\emph{yūndǎo}.) To fight to save a patient? (\emph{qiǎngjiù}.) An earthquake? (\emph{dìzhèn}.) To die? (\textbf{\zh{死}}, \emph{sǐ}.) The life of a living thing? (\emph{shēngmìng}.) The cause of death? (\textbf{\zh{死因}}, \emph{sǐyīn}.) To blow up? (\emph{bàozhà}.) To flee from danger? (\emph{táo}.) In a panic? (\emph{huāng}.) To help up someone who fell? (\emph{fú}.) To break a bone? (\emph{gǔzhé}.) Sure enough? (\textbf{\zh{果然}}, \emph{guǒrán}.)
\end{tcolorbox}
